\chapter{Citações, Notas e Referências}

Este capítulo descreve a operação dos recursos bibliográficos. Os critérios para selecionar fontes, citar com integridade e elaborar referências são tratados na Parte I.

\section{Arquivo bibliográfico}

Cadastre as obras em \texttt{obrigatorio/referencias.bib} e associe o arquivo ao documento principal:

\begin{lstlisting}
\addbibresource{obrigatorio/referencias.bib}
\end{lstlisting}

Cada entrada possui uma chave única, usada pelos comandos de citação. Evite espaços, acentos e chaves duplicadas. Sempre confira no PDF se nomes, título, edição, data, endereço eletrônico e data de acesso foram processados como esperado.

\section{Citações no texto}

\begin{lstlisting}
\citarparenteses{chave}
\citarparenteses[p. 25]{chave}
\citartexto{chave}
\citartexto[p. 25]{chave}

\citardireta[p. 25]{chave}{Texto citado.}
\citardiretatexto[p. 25]{chave}{Texto citado.}
\end{lstlisting}

O argumento entre colchetes é opcional e recebe o localizador, por exemplo \texttt{p. 25}. O argumento seguinte é a chave da entrada bibliográfica.

Para citação direta longa, use:

\begin{lstlisting}
\begin{citacaolonga}[p. 25]{chave}
Texto da citação.
\end{citacaolonga}
\end{lstlisting}

Para citação de citação, prefira informar autoria e ano da original como texto,
usando uma chave apenas para a fonte efetivamente consultada:

\begin{lstlisting}
\parenciteapud{Freire}{1994}[13]{consultada}[25]
\textapud{Freire}{1994}[13]{consultada}[25]
\end{lstlisting}

Os comandos anteriores \texttt{\string\apudbib} e \texttt{\string\apudbibonline}
continuam aceitando duas chaves. Ao usar \texttt{\string\printundfreferences},
a obra original usada somente nessa citação indireta é omitida da lista; uma
citação independente da original na mesma seção bibliográfica a inclui. Essa
filtragem não é aplicada a chamadas diretas a \texttt{\string\printbibliography}.

\section{Notas de rodapé}

\begin{lstlisting}
\notarodape{Texto explicativo da nota.}
\end{lstlisting}

\section{Lista de referências}

Insira \texttt{\string\printundfreferences} após a parte textual. A lista é produzida apenas com as entradas citadas, salvo configuração bibliográfica adicional.

\section{Diagnóstico bibliográfico}

Uma citação com a chave impressa em negrito ou como marcador indefinido geralmente indica chave inexistente ou ciclo de compilação incompleto. Mensagens como \texttt{Please (re)run Biber} pedem a execução do Biber e novas passagens do LuaLaTeX.
